\documentclass[12pt,a4paper]{article}
\usepackage{cmap}
\usepackage[T2A]{fontenc}
\usepackage[utf8]{inputenc}
\usepackage[russian]{babel}
\usepackage[a4paper,margin=2cm]{geometry}
\usepackage{amsmath,amssymb,mathtools}
\usepackage{microtype}
\usepackage[hidelinks]{hyperref}
\input{glyphtounicode}
\pdfgentounicode=1

\setlength{\parindent}{0pt}
\setlength{\parskip}{6pt}

\title{Формулы лабораторных работ}
\author{SMTU-TAU}
\date{\today}

\begin{document}
\maketitle

\section*{1. Основы автоматизации (ОА)}

\subsection*{ЛР1}
\[
\dot{\varphi} = -\frac{\theta}{T_a}\varphi + \frac{1}{T_a}\mu - \frac{1}{T_a}f
\]
\[
\varphi_{i+1} = \varphi_i + \Delta t\,\dot{\varphi}_i
\]

\subsection*{ЛР2}
\[
\mu(t)=\mu_0+\frac{T_r^2}{T_k}\left(\dot{\mu}_0+\frac{\varphi_0}{T_k}\right)\left(1-e^{-\frac{T_k}{T_r^2}t}\right)+\frac{\varphi_0}{T_k}t
\]
\[
\dot{\mu}(t)=\left(\dot{\mu}_0+\frac{\varphi_0}{T_k}\right)e^{-\frac{T_k}{T_r^2}t}+\frac{\varphi_0}{T_k}
\]

\subsection*{ЛР3}
\[
T_r^2\ddot{\mu}+T_k\dot{\mu}+\gamma\mu=\varphi_0,\qquad
D=T_k^2-4T_r^2\gamma
\]
\[
\lambda_{1,2}=\frac{-T_k\pm\sqrt{D}}{2T_r^2}
\]
\[
\mu(t)=e^{\beta t}\bigl(C_1\cos(\omega t)+C_2\sin(\omega t)\bigr)+c,
\quad
\beta=-\frac{T_k}{2T_r^2},
\quad
\omega=\frac{\sqrt{4T_r^2\gamma-T_k^2}}{2T_r^2},
\quad
c=\frac{\varphi_0}{\gamma}
\]

\subsection*{ЛР4}
\[
a\ddot{\mu}+b\dot{\mu}+c\mu=f
\]
\[
a=T_r^2T_a,\quad
b=T_r^2\Theta+T_kT_a,\quad
c=T_k\Theta
\]
\[
\dot{\varphi}=\frac{-\mu-f-\Theta\varphi}{T_a},
\qquad
\varepsilon(t)=-\frac{1}{T_a}\varphi(t)-\mu(t)
\]

\subsection*{ЛР5}
\[
\dot{\mu}=v,\qquad
\dot{v}=\frac{\varphi-T_kv-\gamma\mu}{T_r^2},\qquad
\dot{\varphi}=\frac{-\mu-f-\Theta\varphi}{T_a}
\]
Численный метод: RK4.
\[
\varepsilon(t)=-\frac{1}{T_a}\varphi(t)-\mu(t)
\]

\section*{2. ТАУ Лин}

\subsection*{ЛР1 (дифференцирующее звено с замедлением)}
\[
y(t)=\frac{KX_m}{T}e^{-t/T}
\]
\[
A(\omega)=\frac{K\omega}{\sqrt{T^2\omega^2+1}},
\qquad
\phi(\omega)=\arctan\left(\frac{1}{T\omega}\right)
\]
\[
\Re(\omega)=\frac{KT^2\omega^2}{T^2\omega^2+1},
\qquad
\Im(\omega)=\frac{K\omega}{T^2\omega^2+1}
\]
\[
L(\omega)=20\log_{10}\left(\frac{K\omega}{\sqrt{1+(T\omega)^2}}\right)-20\log_{10}\left(\frac{K}{T}\right),
\quad
\phi_{\mathrm{deg}}(\omega)=90^\circ-\arctan(T\omega)\frac{180^\circ}{\pi}
\]

\subsection*{ЛР2 (апериодическое звено 1-го порядка)}
\[
y(t)=KX_m\left(1-e^{-t/T}\right)
\]
\[
A(\omega)=\frac{K}{\sqrt{T^2\omega^2+1}},
\qquad
\phi(\omega)=-\arctan(T\omega)
\]
\[
\Re(\omega)=\frac{K}{T^2\omega^2+1},
\qquad
\Im(\omega)=-\frac{K\omega}{T^2\omega^2+1}
\]

\subsection*{ЛР3 (апериодическое звено 2-го порядка)}
\[
y(t)=K_1K_2X_m\left(1-\frac{T_1}{T_1-T_2}e^{-t/T_1}+\frac{T_2}{T_1-T_2}e^{-t/T_2}\right)
\]
\[
A(\omega)=\frac{K_1K_2}{\sqrt{(T_1^2\omega^2+1)(T_2^2\omega^2+1)}}
\]
\[
\phi(\omega)=-(\arctan(T_1\omega)+\arctan(T_2\omega))
\]
\[
\Re(\omega)=\frac{K_1K_2(1-T_1T_2\omega^2)}{(1-T_1T_2\omega^2)^2+(T_1+T_2)^2\omega^2},
\quad
\Im(\omega)=-\frac{K_1K_2\omega(T_1+T_2)}{(1-T_1T_2\omega^2)^2+(T_1+T_2)^2\omega^2}
\]

\subsection*{ЛР4 (колебательное звено)}
\[
\omega_n=\frac{\sqrt{1-\xi^2}}{T}
\]
\[
y(t)=KX_m\left[1-e^{-\xi t/T}\left(\cos(\omega_nt)+\frac{\xi}{\sqrt{1-\xi^2}}\sin(\omega_nt)\right)\right]
\]
\[
A(\omega)=\frac{K}{\sqrt{(1-T^2\omega^2)^2+(2\xi T\omega)^2}}
\]
\[
\phi(\omega)=-\operatorname{atan2}(2\xi T\omega,\;1-T^2\omega^2)
\]
\[
\Re(\omega)=\frac{K(1-T^2\omega^2)}{(1-T^2\omega^2)^2+(2\xi T\omega)^2},
\quad
\Im(\omega)=-\frac{K(2\xi T\omega)}{(1-T^2\omega^2)^2+(2\xi T\omega)^2}
\]

\section*{3. ТАУ Нелин}

\subsection*{Базовая динамика (ЛР1-ЛР5)}
\[
x=\varphi-\delta\mu,\qquad
\dot{\varphi}=\frac{1}{T_a}(z-\theta\varphi-\mu),\qquad
\dot{\mu}=\frac{F(x,\dot{x})}{T_s},\qquad
\dot{x}=\dot{\varphi}-\delta\dot{\mu}
\]
Интегратор: явный Эйлер (в части ЛР есть защита от NaN/Inf).

\subsection*{Нелинейности}
\[
F_{\text{2-поз}}(x)=
\begin{cases}
a, & x>0,\\
-a, & x\le 0.
\end{cases}
\]
\[
F_{\text{3-поз}}(x)=
\begin{cases}
a, & x>b,\\
0, & |x|\le b,\\
-a, & x<-b.
\end{cases}
\]
\[
F_{\text{нас}}(x)=
\begin{cases}
Kx, & |x|\le b,\\
a, & x>b,\\
-a, & x<-b.
\end{cases}
\]
\[
F_{\text{нечувст}}(x)=
\begin{cases}
K(x-b), & x>b,\\
0, & |x|\le b,\\
K(x+b), & x<-b.
\end{cases}
\]
\[
F_{\text{нечувст+огр}}(x)=
\begin{cases}
a, & x>c,\\
K(x-b), & b<x\le c,\\
0, & -b\le x\le b,\\
K(x+b), & -c\le x<-b,\\
-a, & x<-c.
\end{cases}
\]
\[
F_{\text{гист}}(x,\dot{x})=
\begin{cases}
a, & \dot{x}>0,\;x>b,\\
-a, & \dot{x}>0,\;x\le b,\\
a, & \dot{x}<0,\;x>-b,\\
-a, & \dot{x}<0,\;x\le -b.
\end{cases}
\]
\[
F_{\text{люфт}}(x,\dot{x})=
\begin{cases}
K(x-b), & \dot{x}>0,\\
K(x+b), & \dot{x}<0.
\end{cases}
\]
\[
F_{\text{нечувст+гист}}(x,\dot{x})=
\begin{cases}
a, & \dot{x}>0,\;x>c,\\
0, & \dot{x}>0,\;-b\le x\le c,\\
-a, & \dot{x}>0,\;x<-b,\\
a, & \dot{x}<0,\;x>b,\\
0, & \dot{x}<0,\;-c\le x\le b,\\
-a, & \dot{x}<0,\;x<-c.
\end{cases}
\]

\section*{4. ТДЗ}

В обозначениях ниже \(s=j\omega\).

\subsection*{ЛР1 (разомкнутая система, вариант 1)}
\[
W_1(s)=\frac{K_1K_2K_3}{(1+T_1s)(1+T_2s)(1+T_3s)}
\]
\[
A(\omega)=\frac{K_1K_2K_3}{\sqrt{(1+T_1^2\omega^2)(1+T_2^2\omega^2)(1+T_3^2\omega^2)}}
\]

\subsection*{ЛР2 (разомкнутая система, вариант 2)}
\[
W_2(s)=\frac{K_1K_2}{(1+T_1s)\left(1+2\xi T_2s+T_2^2s^2\right)}
\]
\[
A(\omega)=\frac{K_1K_2}{\sqrt{(1+T_1^2\omega^2)\left((1-T_2^2\omega^2)^2+(2\xi T_2\omega)^2\right)}}
\]

\subsection*{ЛР3 (замкнутая система, вариант 1)}
\[
H_3(s)=\frac{K_1K_2(1+T_3s)}
{T_1T_2T_3s^3+(T_1T_2+T_1T_3+T_2T_3)s^2+(T_1+T_2+T_3)s+(1+K_1K_2K_3)}
\]
\[
D_3(s)=T_1T_2T_3s^3+(T_1T_2+T_1T_3+T_2T_3)s^2+(T_1+T_2+T_3)s+(1+K_1K_2K_3)
\]

\subsection*{ЛР4 (замкнутая система, вариант 2)}
\[
H_4(s)=\frac{K_1K_2(1+T_2s)}
{T_1^2T_2s^3+(2\xi T_1T_2+T_1^2)s^2+(T_2+2\xi T_1)s+(1+K_1K_2)}
\]
\[
D_4(s)=T_1^2T_2s^3+(2\xi T_1T_2+T_1^2)s^2+(T_2+2\xi T_1)s+(1+K_1K_2)
\]

\subsection*{ЛР5 (замкнутая система, вариант 3)}
\[
H_5(s)=\frac{K_1(1+T_2s)(1+T_3s)}
{T_1T_2T_3s^3+(T_1T_2+T_1T_3+T_2T_3)s^2+(T_1+T_2+T_3)s+(1+K_1K_2K_3)}
\]
\[
D_5(s)=T_1T_2T_3s^3+(T_1T_2+T_1T_3+T_2T_3)s^2+(T_1+T_2+T_3)s+(1+K_1K_2K_3)
\]

\subsection*{ЛР6 (замкнутая система, вариант 4)}
\[
H_6(s)=\frac{K_1(T_2^2s^2+2\xi T_2s+1)}
{T_1T_2^2s^3+(2\xi T_1T_2+T_2^2)s^2+(T_1+2\xi T_2)s+(1+K_1K_2)}
\]
\[
D_6(s)=T_1T_2^2s^3+(2\xi T_1T_2+T_2^2)s^2+(T_1+2\xi T_2)s+(1+K_1K_2)
\]

\end{document}
